\section{Reproduction Setup}
\label{sec:setup}

\subsection{Starting Code}
We start from the authors' bug collection~\cite{gobugsrepo}
(\url{https://github.com/system-pclub/go-concurrency-bugs}).
\todo{Commit hash used, and the list of files we added or modified.}

\subsection{Environment}
\begin{table}[h]
  \centering
  \caption{Reproduction environment.}
  \label{tab:env}
  \begin{tabular}{ll}
    \toprule
    Item & Value \\
    \midrule
    CPU / RAM         & \todo{} \\
    Host OS / kernel  & \todo{} \\
    Container image   & Ubuntu 18.04 (Docker) \\
    Go toolchain      & Go 1.10.8 (linux/amd64) \\
    Docker version    & \todo{} \\
    \bottomrule
  \end{tabular}
\end{table}

\subsection{Workloads}
\todo{The bug cases we run (from \texttt{bugs/channel/} and \texttt{bugs/shared\_memory/}),
how they are obtained, and how each run is triggered and judged (e.g., detector output, timeout).}

\subsection{Deviations}
\begin{table}[h]
  \centering
  \caption{Deviations from the paper's setup.}
  \label{tab:deviations}
  \begin{tabular}{p{3cm}p{3.5cm}p{3.5cm}p{4cm}}
    \toprule
    Item & Original setting & Setting used & Reason \\
    \midrule
    \todo{} & & & \\
    \bottomrule
  \end{tabular}
\end{table}
